\chapter{Introduzione}

% Nella stesura  della tesi, in quasi tutti i casi tre parti devono emergere con chiarezza:
% \begin{itemize}
% 	\item analisi dello stato dell’arte e definizione del problema
% 	\item studio e/o sviluppo di modelli per la risoluzione del problema
% 	\item risultati teorici o sperimentali ottenuti
% \end{itemize} 

% La scelta di come strutturare l'elaborato finale non è univoca, pertanto si suggerisce di iniziare creando un indice organizzato in maniera gerarchica suddiviso in capitoli, sezioni, sottosezioni (non eccedere con la profondità, si rischia di frammentare troppo il testo).

% Un esempio di 'albero' che prevede annidamenti fino alle sottosezioni è il seguente (puramente indicativo):

\section{Il problema affrontato}

La ricostruzione tridimensionale della forma e della posa delle mani a partire da input visivi rappresenta una delle sfide più complesse e attuali nell'ambito della Computer Vision.
Le mani e il corpo operano in costante sinergia per comunicare e risolvere compiti complessi; pertanto, una comprensione completa delle azioni, delle emozioni e delle intenzioni umane non è possibile senza un'analisi congiunta di entrambi.
Il tracciamento del movimento delle mani trova molteplici applicazioni nell'ambito dell'interazione uomo-macchina (Human-Computer Interaction, HCI) per lo sviluppo di interfacce più naturali e intuitive, nonché in ambito di telemedicina, robotica ed animazione.

Sorprendentemente, la modellazione 3D e il tracciamento del corpo e delle mani sono stati storicamente affrontati nella ricerca come due domini separati \cite{MANO:SIGGRAPHASIA:2017,EROL200752}.

La stima della posa del corpo umano è stata in realtà un campo di ricerca a lungo molto più attivo.
Questo ha portato allo sviluppo di tecniche di tracciamento che hanno poi potuto essere riutilizate negli algoritmi di tracciamento delle mani.

Esistono tuttavia notevoli differenze tra i due ambienti operativi:
per esempio, la presenza dei vestiti sul corpo aumenta la complessità di stima della posa, ma, al contempo, la presenza di colori aiuta la segmentazione e l'estrazione di caratteristiche basate su colore e texture;
in particolare nei confronti dell'uniformità di colore delle mani.
Un altro esmepio può essere la possibilità di stimare la posa del corpo umano parte per parte o in maniera gerarchica (per prima la testa, poi torso, gambe, etc.), così da suddividere il problema in sottoproblemi con una complessità dimensionale inferiore.
Nel caso della mano, un approccio gerarchio rimarrebbe limitato a due sole parti: il palmo e le dita \cite{EROL200752}.

Dal punto di vista tecnico, il problema dell'inferenza della posa e della ricostruzione della mesh delle mani è intrinsecamente difficile a causa di molteplici fattori.

\begin{itemize}
    \item La mano umana gode di una notevole mobilità; essa costituisce infatti un sistema articolato con oltre 20 gradi di libertà (DOF: \emph{degree of freedom}).
    Sebbene una mano reale non goda di 20 gdl completi a causa dell'interdipendenza tra dita, studi hanno mostrato come sia impossibile modellarne una con meno di 6 gdl.
    
    \item Nelle immagini o nelle scansioni 3D a corpo intero, le mani occupano una porzione minima dello spazio visivo o del volume di acquisizione.
    Il tracciamento visivo è inoltre complicato dalla forte somiglianza strutturale e cromatica tra le singole dita.
    Queste caratteristiche rendono le dita estremamente difficili da risolvere, generando spesso rumore o l'effetto di fusione dei vertici delle dita (\emph{webbing}).
    In aggiunta a ciò, la naturale geometria delle mano insieme alla frequente interazione con oggetti causa continue occlusioni ed auto-occlusioni, rendendo l'inferenza della posa un problema altamente ambiguo.

    \item La cinematica tipica delle mani è caratterizzata da movimenti estremamente rapidi, che possono raggiungere i~$5$\,m/s per le traslazioni e i~$300$\,°/s per le rotazioni del polso; ciò genera frequenti fenomeni di \emph{motion blur} (sfocatura da movimento) nelle acquisizioni, degradando la qualità dell'immagine.

    \item Anche per una singola sequenza di immagini, un sistema di computer vision deve essere in grado di elaborare un'enorme quantità di dati, spesso rispettando, al contempo, stringenti requisiti di latenza.
    Gli algoritmi di tracciamento devono perciò essere altamente efficienti e spesso richiedono hardware specifico ottimizzato per il calcolo parallelo per l'esecuzione.

    \item Per garantire un'ampia diffusione, il sistema deve poter operare in ambienti non controllati, con illuminazione variabile e sfondi complessi.
    Tuttavia, segmentare un oggetto rigido in un ambiente arbitrario costituisce già di per sè un compito non banale.
\end{itemize}

Per far fronte a queste problematiche, in passato molti algoritmi per la stima della posa della mano ponevano stringenti vincoli all'utente o sull'ambiente di acquisizione.
Ad esempio, poteva essere richiesto uno sfondo uniforme, in cui la mano fosse l'unico oggetto mobile e colorato; di evitare rapidi movimenti della mano; di mantenere il palmo parallelo al piano immagine; di campionare con una frequenza molto elevata.
% TODO: verificare. Erol 2.1
Questi vincoli erano tutti volti a ridurre i gradi di libertà della mano e semplificare il problema.

Per queste ragioni i metodi migliori per l'acquisizione del movimento della mano si sono basati a lungo sull'utilizzo di dispositivi elettromeccanici o magnetici da indossare, comunemente noti come \emph{data gloves}.
Sebbene tali dispositivi forniscano misurazioni complete in tempo reale, essi presentano notevoli limitazioni pratiche, in quanto tali sensori risultano onerosi dal punto di vista economico, ostacolano la naturalezza del movimento e richiedono prolungate procedure di calibrazione e configurazione per ottenere stime precise.
Inoltre, i sistemi di tracciamento basati sull'applicazione di marcatori (marker-based) si rivelano spesso intrusivi, in quanto i marcatori stessi possono alterare la forma della mano e limitarne l'ampiezza dei movimenti.
È stato inoltre osservato come i marcatori possono staccarsi durante l'uso e sono anch'essi soggetti ad occlusioni, necessitando pertanto di un numero elevato di telecamere per garantire una corretta risoluzione in un volume di cattura esteso.

L'approccio basato sulla visione artificiale (Computer Vision) offre il potenziale per superare tali barriere, proponendo soluzioni non a contatto, prive di ingombri fisici e conseguentemente più naturali \cite{EROL200752}.


\section{Il problema della modellazione della mano}


\subsection{Il modello cinematico della mano}

Il modello cinematico della mano si basa sulla rappresentazione dello scheletro come un sistema di corpi rigidi interconnessi da articolazioni dotate di uno o più gradi di libertà (Degrees of Freedom, DOF) rotazionali.
Dal punto di vista anatomico e cinematico, il modello si articola secondo le seguenti caratteristiche:

\begin{itemize}
    \item La struttura ossea è composta da 27 ossa complessive, di cui 8 localizzate nel polso e le restanti 19 strutturate a formare il palmo e le dita.
    
    \item Le articolazioni interfalangee (Interphalangeal, IP), che collegano i segmenti interni di ciascun dito, sono nove in totale e vengono modellate in modo accurato con un singolo grado di libertà adibito ai movimenti di flessione ed estensione.
    
    \item Le articolazioni metacarpo-falangee (Metacarpophalangeal, MCP), responsabili dell'unione tra le dita e il palmo, sono tipicamente descritte in letteratura come giunti a sella con due gradi di libertà.
    Tali gradi consentono i movimenti di flessione ed estensione, nonché l'abduzione e l'adduzione (l'allontanamento reciproco delle dita) sul piano definito dal palmo.
    
    \item Le articolazioni carpo-metacarpali (Carpometacarpal, CMC) collegano le ossa metacarpali al polso.
    Le articolazioni CMC del dito indice e del dito medio sono considerate statiche. 
    Sebbene le articolazioni CMC dell'anulare e del mignolo possiedano una mobilità limitata associata alla curvatura del palmo, tale movimento è frequentemente trascurato, portando all'assunzione di un palmo completamente rigido.
    
    \item L'articolazione CMC del pollice, denominata trapeziometacarpale (TM), rappresenta il giunto di più complessa modellazione, in quanto dotata di due assi di rotazione non ortogonali e non secanti. 
    Il modello a sella con due gradi di libertà risulta limitante per questo giunto; per superare tale restrizione è stata proposta in alcuni studi l'estensione a un giunto sferico con tre gradi di libertà.
\end{itemize}

La combinazione dei gradi di libertà angolari delle dita, definita \emph{configurazione locale}, e dei sei gradi di libertà associati a un sistema di riferimento posizionato nel polso, definita \emph{configurazione globale}, costituisce il vettore di configurazione che rappresenta la posa completa della mano.

Un modello cinematico standard e ampiamente utilizzato nella ricerca prevede complessivamente 27 gradi di libertà. 
In questo modello, assumendo le articolazioni CMC come fisse, il palmo viene modellato in maniera non del tutto realistica come un unico corpo rigido. 
Le dita sono invece schematizzate come catene cinematiche seriali planari, agganciate al palmo mediante punti di ancoraggio coincidenti con le articolazioni MCP.

Poiché l'assunzione di planarità per il movimento delle dita non è generalmente valida nella realtà, sono state introdotte diverse varianti per affinare il modello cinematico.
Tra le soluzioni adottate figurano:

\begin{itemize}
    \item L'aggiunta di un movimento di torsione supplementare alle articolazioni MCP.
    \item L'inserimento di un movimento di torsione attorno all'asse osseo come funzione lineare degli angoli di abduzione e flessione.
    \item L'introduzione di un grado di libertà per la flessione e l'estensione nelle articolazioni CMC.
    \item L'impiego di giunti sferici dedicati specificamente all'articolazione TM del pollice.
\end{itemize}

% inserire immagine con RX mano dal paper e generare schema della struttura


\subsection{La modellazione del movimento naturale}

Sebbene la mano sia in grado di assumere innumerevoli conformazioni, il movimento attivo risulta in realtà fortemente vincolato dalla biologia.
Vincoli che non sono riprodotti dal modello cinematico.
Pertanto, affinchè gli algoritmi di stima operino in modo efficiente ed evitino configurazioni impossibili, è indispensabile introdurre dei vincoli che circoscrivano lo spazio di ricerca alle sole pose fisiologicamente plausibili.

Un primo approccio alla risoluzione di tale problematica prevede la definizione di vincoli analitici derivati dagli studi sulla biomeccanica.
Tali formulazioni si suddividono in vincoli\emph{statici}, che modellano l'escursione di ogni parametro, e \emph{dinamici}, i quali modellano le dipendenze meccaniche tra giunti adiacenti.
Un classico esempio è il vincolo esistente tra le articolazioni distale e prossimale, espresso dalla relazione $\theta_{DIP} = \frac{2}{3}\theta_{PIP}$.

Tuttavia, l'articolata struttura della mano insieme alle complesse interazioni anatomiche impediscono di esprimere tutti i vincoli in forma chiusa.
Per far fronte a queste complessità gli studi si sono spostati verso approcci basati sull'apprendimento automtico, sfruttando dati \emph{ground truth} acquisiti tramite sensori indossabili come i data gloves.
L'assunto fondamentale di questi approcci è che le pose cinematicamente realizzabili risiedano in una varietà topologica (manifold) a dimensionalità notevolmente inferiore rispetto allo spazio di partenza.
Applicando tecniche di riduzione dimensionale, come l'Analisi delle Componenti Principali (PCA) su dati articolari, è stato possibile approssimare la posa della mano in uno spazio a sette dimensioni \cite{lin2000modeling}.
In alternativa, i dati campionati possono essere impiegati in forma diretta per generare ampi database di template sintetici, utili in fase di tracciamento per ricnoscere la mano in tutte le pose.

Infine, una corretta modellazione del movimento naturale richiede non solo la restrizione spaziale delle pose ammissibili, ma anche l'apprendimento della dinamica temporale del moto. Per modellare le complesse dinamiche non lineari, sono stati proposti metodi basati sull'analisi dinamica agli autovalori (eigen-dynamic analysis, EDA).
Tali metodi impiegano inizialmente la PCA per la riduzione dimensionale e, successivamente, modellano i movimenti indipendenti delle dita tramite sistemi lineari di basso ordine, combinandoli infine in un sistema dinamico lineare stocastico di ordine superiore.Metodologie alternative prevedono la schematizzazione dello spazio delle configurazioni sotto forma di albero di ricerca, costruito mediante tecniche di clustering gerarchico o partizionamento dell'autospazio. 


\subsection{Modellazione della forma}



\section{Stato dell'arte}

% TODO inserire sec 2.2 taxonomy dal paper di Erol
% parlo dei modelli già esistenti: MANO, SMPL, Hamer

\section{Soluzioni in letteratura}
